\documentclass[11pt]{amsart}

\usepackage{amsmath,amssymb,mathrsfs}
\usepackage[centering,text={15.5cm,24cm}]{geometry}
%\usepackage[centering,paperwidth=18cm,paperheight=18.5cm,
%text={16.5cm,16.5cm},top=1.25cm]{geometry} % laptop
%\usepackage[centering,paperwidth=23cm,paperheight=30cm,
%text={20cm,28cm},top=1.25cm]{geometry} % desktop
\usepackage{xcolor}
\usepackage{hyperref}

\hypersetup{
  colorlinks=true,
  linkcolor=blue,
  citecolor=violet,
  urlcolor=violet
}

\renewcommand{\baselinestretch}{1.2}
\allowdisplaybreaks[1]

\numberwithin{equation}{section}

% Compact section headings
\makeatletter
\renewcommand{\section}{%
  \@startsection{section}{1}%
    \z@
    {.5\linespacing\@plus.7\linespacing}%
    {.35\linespacing}%
    {\normalfont\bfseries}}
\makeatother

\theoremstyle{plain}
\newtheorem{theorem}{Theorem}[section]
\newtheorem{proposition}[theorem]{Proposition}
\newtheorem{lemma}[theorem]{Lemma}
\newtheorem*{maintheorem}{Main Theorem}

\theoremstyle{definition}
\newtheorem*{question}{Question}

\newcommand{\CC}{\mathbb{C}}
\newcommand{\NN}{\mathbb{N}}
\newcommand{\QQ}{\mathbb{Q}}
\newcommand{\RR}{\mathbb{R}}

\newcommand{\cN}{\mathcal{N}}
\newcommand{\scrL}{\mathscr{L}}
\newcommand{\scrN}{\mathscr{N}}
\newcommand{\delbar}{\bar\partial}

\title{Almost complex structures making a symplectic snc divisor
pseudo-holomorphic}

\author{Bernd Siebert}

\date{\today}

\begin{document}
\begin{center}
{\large Bernd Siebert}
\end{center}
\begin{maintheorem}
Let $(M,\omega)$ be a closed symplectic manifold of real dimension $2n$ and
let $D=D_1\cup\cdots\cup D_s\subset M$ be a symplectic simple normal
crossings divisor. Then there exist a path $\omega_t$, $t\in[0,1]$, of
cohomologous symplectic forms with $\omega_0=\omega$, along which $D$ remains
a symplectic simple normal crossings divisor, and an $\omega_1$-tame almost
complex structure $J$ such that every $p\in D$ has a neighborhood $U$ with
smooth complex coordinates $(z_1,\ldots,z_n)$ centered at $p$ for which, for
some $1\le r\le n$, the functions $z_1,\ldots,z_r$ are $J$-holomorphic and
\[
D\cap U=\{z_1\cdots z_r=0\}.
\]
\end{maintheorem}

The proof combines the regularization theorem of Tehrani--McLean--Zinger
\cite{TMZ1} with a relative simultaneous version of Donaldson's theorem
\cite{Donaldson}, using Auroux's vector-bundle generalization \cite{Auroux}, and
then argues by pseudo-holomorphic structures on normal bundles constructed
inductively by pseudo-holomorphic divisors. We begin by discussing the latter
two notions.

%------------------------------------------------------------------------------

\section{Pseudo-holomorphic line bundles and divisors}
\label{SS: Pshol line bundles}

Let $(X,J)$ be an almost complex manifold. A $(0,1)$-connection on a complex
vector bundle $E\to X$ is a $\CC$-linear map $\delbar_E:C^\infty(E)\to
C^\infty(E\otimes\Lambda_X^{0,1})$ satisfying
$\delbar_E(fs)=s\otimes\delbar_Jf+f\delbar_Es$. It is a \emph{$J$-holomorphic
structure} if $E$ locally has frames $e_1,\ldots,e_m$ with $\delbar_Ee_i=0$.
Such a bundle will also be called \emph{pseudo-holomorphic}. A morphism
$\varphi:E\to F$ of pseudo-holomorphic vector bundles is a smooth complex-linear
bundle map satisfying $\delbar_F(\varphi s) =
(\varphi\otimes\mathrm{id})\delbar_Es$ for every smooth section $s$ of $E$. An
isomorphism is an invertible morphism. Equivalently, pseudo-holomorphic bundles
have $J$-holomorphic transition functions, and their morphisms are represented
in pseudo-holomorphic frames by matrices of $J$-holomorphic functions. Direct
sums, tensor products and duals have the induced structures.

A pseudo-holomorphic vector bundle $E$ has a natural almost complex structure
on its total space. Indeed, in a pseudo-holomorphic frame let
$z_1,\ldots,z_m$ be the linear fiber coordinates and declare the
$(1,0)$-forms to be generated by the pullbacks of $\Lambda_X^{1,0}$ and
$dz_1,\ldots,dz_m$. A pseudo-holomorphic change of frame preserves this
subbundle. The resulting almost complex structure is the unique one which
induces the given complex structure on every fiber and for which the
projection $E\to X$ and all pseudo-holomorphic sections are
pseudo-holomorphic.

A smooth codimension-two submanifold $Z\subset X$ is a
\emph{$J$-holomorphic divisor} if it is locally the transverse zero locus of
a $J$-holomorphic function. A union $Z_1\cup\cdots\cup Z_s$ is a
\emph{$J$-holomorphic snc divisor} if the branches through any point have
$J$-holomorphic defining functions $f_1,\ldots,f_r$ such that
$(f_1,\ldots,f_r)$ is a submersion. The smooth constant-rank theorem
completes these functions to complex-valued smooth coordinates, so this is
the local condition in the Main Theorem.

The following elementary fact is responsible for the rigidity of smooth
pseudo-holomorphic divisors.

\begin{lemma}
\label{Lem: f/g is J-hol for smooth D}
Let $f,g:U\to\CC$ be $J$-holomorphic functions having the same transverse,
positively oriented zero locus $Z$. Then $f/g$ extends uniquely to a
nowhere-vanishing $J$-holomorphic function on $U$.
\end{lemma}

\begin{proof}
The assertion is local near a point $p\in Z$. Since $g$ is a submersion, we may
take $z_1=g$ as one of a system of smooth complex-valued coordinates centered at
$p$, so $Z=\{z_1=0\}$ and $\delbar_Jz_1=0$. Write $z_1=x_1+iy_1$. Integrating
the radial derivative of $f(tx_1,ty_1,z')$ from $t=0$ to $t=1$, and using
$f|_Z=0$, gives $f=az_1+b\bar z_1$ for smooth functions $a,b$. Hence
\[
0=\delbar_Jf
=z_1\delbar_Ja+b\delbar_J\bar z_1+\bar z_1\delbar_Jb.
\]
Restricting to $Z$ gives $b|_Z=0$. Thus $b=O(|z_1|)$ and $f/z_1=a+b\bar z_1/z_1$
extends continuously across $Z$, with value $a|_Z$. This value is nonzero by
transversality and the orientation assumption. The quotient is locally bounded
and $J$-holomorphic on $U\setminus Z$. Since $Z$ has real codimension two,
\cite[Thm.\,4.1(b)]{HarveyPolking}, applied to $\delbar_J$, shows that it is a
weak solution across $Z$; elliptic regularity then makes it smooth and
$J$-holomorphic.
\end{proof}

\begin{proposition}
\label{Prop: O(D)}
For every smooth $J$-holomorphic divisor $Z\subset X$ there is a
pseudo-holo\-mor\-phic line bundle $\scrL_X(Z)$ with a pseudo-holomorphic
section $s_Z$ having transverse zero locus $Z$. The pair $(\scrL_X(Z),s_Z)$ is
unique up to a unique pseudo-holomorphic isomorphism preserving the section.
Conversely, every transverse pseudo-holomorphic section of a pseudo-holomorphic
line bundle defines a smooth $J$-holomorphic divisor.
\end{proposition}

\begin{proof}
Choose local $J$-holomorphic defining functions $f_\alpha$ for $Z$, taking
$f_\alpha=1$ on open sets disjoint from $Z$. By
Lemma~\ref{Lem: f/g is J-hol for smooth D}, the quotients
$f_\alpha/f_\beta$ extend to nowhere-vanishing $J$-holomorphic functions.
They are the transition functions of $\scrL_X(Z)$, and the $f_\alpha$ define
its canonical section $s_Z$. If $(L,s)$ is another such pair, the local
quotients of $s$ by $s_Z$ extend by the same lemma and give the unique
pseudo-holomorphic isomorphism taking $s_Z$ to $s$. The converse follows by
writing a section in a pseudo-holomorphic frame.
\end{proof}

For an integral combination $E=\sum_i a_iZ_i$ of smooth $J$-holomorphic
divisors, set
$\scrL_X(E)=\bigotimes_i\scrL_X(Z_i)^{\otimes a_i}$, with negative tensor
powers interpreted as powers of the dual. In particular, for an snc divisor
$Z=\bigcup_iZ_i$, the notation $\scrL_X(Z)$ means
$\bigotimes_i\scrL_X(Z_i)$. Thus Proposition~\ref{Prop: O(D)} and
Lemma~\ref{Lem: f/g is J-hol for smooth D} are applied separately to the
smooth branches; no assertion about arbitrary singular zero loci is being
made.

The normal line bundle $N_{Z|X}$ of a smooth $J$-holomorphic divisor also has a
natural pseudo-holomorphic structure. A local defining function $f$ trivializes
it by $df:N_{Z|X}\to Z\times\CC$. If $g$ is another defining function and
$f=hg$, then $h$ is nowhere-vanishing and $J$-holomorphic by Lemma~\ref{Lem: f/g
is J-hol for smooth D}, while $df|_{N_{Z|X}}=h\,dg|_{N_{Z|X}}$ along $Z$. These
trivializations define a pseudo-holomorphic line bundle, denoted $\scrN_{Z|X}$.
In the local frame $e_f$ with $s_Z=f e_f$, the normal derivative of $s_Z$
sends $v\in N_{Z|X}$ to $df(v)e_f$. This gives the canonical
pseudo-holomorphic identification
\begin{equation}
\label{Eqn: divisor-normal-bundle}
\scrL_X(Z)|_Z\simeq\scrN_{Z|X}.
\end{equation}
We also use restrictions, tensor products and duals of this isomorphism.

%------------------------------------------------------------------------------

\section{A relative Donaldson theorem}
\label{SS: Donaldson theorem}

We need the following relative version of Auroux's vector-bundle
generalization of Donaldson's theorem \cite[Thm.\,2]{Auroux}. An
$\omega$-regularization is understood in the sense of
\cite[Def.\,2.12]{TMZ1}.

\begin{theorem}
\label{Thm: strong Bertini} Let $(X,\omega)$ be a closed symplectic manifold
with $[\omega]$ integral, and let $D\subset X$ be a symplectic snc divisor
equipped with an $\omega$-regularization. Let $L_\omega$ be a complex line
bundle with $c_1(L_\omega)=[\omega]$, and let $L_1,\ldots,L_q$ be complex line
bundles on $X$. Then, for every sufficiently large integer $N$, there are
smooth, transverse sections $s_i$ of $L_i\otimes L_\omega^{\otimes N}$ such
that, for $H_i=s_i^{-1}(0)$, the union $D\cup H_1\cup\cdots\cup H_q$ is a
symplectic snc divisor.
\end{theorem}

\begin{proof}
For a set $I$ of components of $D$, write $D_I$ for their intersection and put
$D_\emptyset=X$. We refer to the nonempty intersections $D_I$, which need not be
connected, as the (closed) strata of $D$. Choose an $\omega$-compatible almost
complex structure $J$ adapted to the regularization, in the sense that all $D_I$
are $J$-invariant and the normal models split along them. Fix the metric
$g=\omega(\,\cdot\,,J\,\cdot\,)$, a Hermitian connection $\nabla_L$ on
$L_\omega$ with curvature $-2\pi i\omega$, and Hermitian metrics and connections
on the $L_i$.

We recall the terminology and the part of Auroux's modified proof of
\cite{Donaldson} used below. For a fixed Hermitian vector bundle $E$, a sequence
$s_k$ of sections of $E\otimes L_\omega^k$ is \emph{asymptotically
$J$-holomorphic} if
\begin{equation}
\label{Eqn: AH estimates}
|s_k|=O(1),\quad |\nabla s_k|=O(k^{1/2}),\quad
|\delbar_Js_k|=O(1),\quad |\nabla^2s_k|=O(k),\quad
|\nabla\delbar_Js_k|=O(k^{1/2}).
\end{equation}
Here $\nabla$ and the norm are induced by $\nabla_L$ and a fixed connection on
$E$ and Riemannian metric on $X$. The sequence is \emph{$\eta$-transverse} on a
$J$-invariant submanifold $Y$ for some $\eta>0$ if, whenever $x\in Y$ and
$|s_k(x)|<\eta$, the restriction $\nabla s_k(x)|_{T_xY}$ has a right inverse of
norm at most $\eta^{-1}k^{-1/2}$. Thus its zero locus in $Y$ is transverse, with
a quantitative lower bound on the normal derivative. We regard the zero section
of the zero bundle as $\eta$-transverse.

Auroux's vector-bundle theorem \cite[Thm.\,2]{Auroux} produces asymptotically
holomorphic sequences of sections of $E\otimes L_\omega^k$ that are uniformly
transverse, meaning $\eta$-transverse for some $\eta>0$. The local
transversality argument proceeds, in local trivializations of $E$, one summand
at a time. If $E=E'\oplus\CC$, let $w_k$ and $\sigma_k$ be asymptotically
holomorphic sequences of sections of $E'\otimes L_\omega^k$ and $L_\omega^k$,
respectively, with $w_k$ uniformly transverse. Perturbing $\sigma_k$ by sums of
localized Gaussian sections makes the restriction $\sigma_k|_{W_k}$ to
$W_k=w_k^{-1}(0)$ uniformly transverse, and hence makes $(w_k,\sigma_k)$
uniformly transverse. The perturbations may be taken arbitrarily small,
uniformly in $k$, in the scaled norms
\begin{equation}
\label{Eqn: scaled C1 norm}
\|u\|_{C^1_k}=\sup_X\bigl(|u|+k^{-1/2}|\nabla u|\bigr).
\end{equation}
This flexibility is what allows the componentwise construction to be repeated
while retaining earlier transversality conditions.

We recall the local sections entering this construction. In Darboux coordinates
$z=(z_1,\ldots,z_n)$ centered at $x$, and in the radial trivialization of
$L_\omega$, the model section of $L_\omega^k$ is the Gaussian $\exp(-k|z|^2/4)$.
Multiplying by a cutoff which equals one on the ball of radius $\frac12k^{-1/3}$
and vanishes outside the ball of radius $k^{-1/3}$ gives a global peak section
$\sigma_{k,x}$. On the cutoff region the Gaussian is $O(e^{-c k^{1/3}})$ for
some $c>0$, so the cutoff error is negligible to every fixed order. The
resulting peak sections satisfy \eqref{Eqn: AH estimates}. Their centers are
chosen in a $k^{-1/2}$-dense set $S=S_1\cup\cdots\cup S_\ell$, where $\ell$ is
independent of $k$ and the points in each $S_j$ are uniformly separated at scale
$k^{-1/2}$. Following \cite{Donaldson}, Auroux perturbs successively over the
$S_j$. Gaussian decay controls interference within each step, while decreasing
perturbation sizes preserve the transversality obtained in earlier steps.

We need two modifications of this construction, both already implicit in the
proof of the vector-bundle theorem. First, suppose that the new component is to
take values in $L_i\otimes L_\omega^k$ rather than $L_\omega^k$. Near each
center choose a unitary frame $e_{i,x}$ of $L_i$ and replace $\sigma_{k,x}$ by
$e_{i,x}\otimes\sigma_{k,x}$. Since $L_i$ is fixed and $X$ is compact, these
frames can be chosen so that the connection matrix of $L_i$ and its first
derivative are uniformly bounded. The additional terms obtained by
differentiating $e_{i,x}$ are therefore of lower order than the corresponding
derivatives of $\sigma_{k,x}$. Hence all estimates in \eqref{Eqn: AH estimates},
as well as the Gaussian summation estimates, remain unchanged. This is the same
local argument used by Auroux for a general fixed vector bundle.

Second, the construction can be made relative to a $J$-invariant symplectic
submanifold $Y$. Suppose that $w_k$ is uniformly transverse on both $X$ and
$Y$. First run the componentwise construction on $Y$, using centers in $Y$.
Choose approximately holomorphic Darboux coordinates on $Y$ at each center
and extend them to Darboux coordinates on $X$ adapted to $Y$. In these
coordinates, $Y$ is a complex coordinate subspace, and the ambient Gaussian
peak section restricts to the Gaussian peak section used on $Y$. Thus the
perturbation constructed on $Y$ is the restriction of an ambient
asymptotically holomorphic perturbation. It makes the new component
transverse on $Y\cap w_k^{-1}(0)$, and hence makes the enlarged vector-bundle
section transverse on $Y$.

One then runs the ambient componentwise construction. Its perturbation may be
chosen so small in the norm \eqref{Eqn: scaled C1 norm} that the
transversality already obtained on $Y$ is preserved. Indeed, if a section is
$\eta$-transverse, a sufficiently small $C^1_k$-perturbation is
$\eta/2$-transverse. The same argument can be repeated for finitely many
$J$-invariant submanifolds, taking each new perturbation small enough to
preserve all conditions obtained before it. We shall apply this observation
to the finite collection of strata of $D$.

For $A\subseteq\{1,\ldots,q\}$ set
$E_A=\bigoplus_{a\in A}L_a$ and, given sections $s_{a,k}$ of
$L_a\otimes L_\omega^k$, write
$s_{A,k}=\bigoplus_{a\in A}s_{a,k}$. We now construct the sections
successively so that
\begin{equation}
\label{Eqn: all subcollections transverse}
s_{A,k}|_{D_I}\ \text{is uniformly transverse for every }D_I
\text{ and every nonempty }A.
\end{equation}

Suppose $s_{1,k},\ldots,s_{m-1,k}$ have been constructed so that
$s_{A,k}|_{D_I}$ is uniformly transverse for every stratum $D_I$ and every
nonempty $A\subseteq\{1,\ldots,m-1\}$. The base case $m=1$ gives
$A=\emptyset$. Begin the construction of the $m$-th component with the zero
section of $L_m\otimes L_\omega^k$. For every stratum $D_I$ and every
$A\subseteq\{1,\ldots,m-1\}$, apply the relative componentwise construction to
make $s_{A\cup\{m\},k}|_{D_I}$ uniformly transverse. When $A=\emptyset$, this
means making $s_{m,k}|_{D_I}$ transverse.

After each such application choose the perturbation small enough to preserve all
previously established conditions. There are finitely many pairs $(I,A)$, so the
process terminates with a positive uniform transversality constant. The total
perturbation is a finite sum of asymptotically holomorphic perturbations and
therefore still satisfies \eqref{Eqn: AH estimates}. It establishes \eqref{Eqn:
all subcollections transverse} for every subset of $\{1,\ldots,m\}$; conditions
involving only the first $m-1$ components have not changed. Induction on $m$
gives sections $s_{1,k},\ldots,s_{q,k}$ satisfying \eqref{Eqn: all
subcollections transverse}.

We finally verify that the zero loci have the required symplectic and
orientation properties. Fix a stratum $D_I$, a nonempty subset $A$, and a
point
\[
p\in Z_{I,A,k}:=D_I\cap\bigcap_{a\in A}s_{a,k}^{-1}(0).
\]
The rescaled derivative
\[
B_k=k^{-1/2}\nabla(s_{A,k}|_{D_I})(p):
T_pD_I\longrightarrow(E_A\otimes L_\omega^k)_p
\]
is surjective and has a right inverse of uniformly bounded norm. Its
complex-antilinear part satisfies $|B_k^{0,1}|=O(k^{-1/2})$ by
\eqref{Eqn: AH estimates}. Thus the complex-linear part $B_k^{1,0}$ is also
surjective for sufficiently large $k$: a right inverse for $B_k$ gives one
for $B_k^{1,0}$ by the usual Neumann-series argument. The same estimate shows
that $\ker B_k$ is $O(k^{-1/2})$-close in the Grassmannian to the complex
subspace $\ker B_k^{1,0}$. Since the complex Grassmannian is contained in the
open set of symplectic planes, $\ker B_k$ is symplectic for large $k$.
Moreover,
\[
T_pZ_{I,A,k}=\ker B_k,
\]
so $Z_{I,A,k}$ is a symplectic submanifold of the expected codimension.

It remains to compare orientations. For $t\in[0,1]$ consider the
surjective maps $B_{k,t}=B_k^{1,0}+tB_k^{0,1}$. For large $k$ all their
kernels remain close to $\ker B_k^{1,0}$ and hence remain symplectic. At $t=0$, 
the exact sequence of complex vector spaces
\[
0\longrightarrow\ker B_{k,0}\longrightarrow T_pD_I
\xrightarrow{\,B_{k,0}\,}
(E_A\otimes L_\omega^k)_p\longrightarrow0
\]
induces the complex, hence symplectic, orientation on $\ker B_{k,0}$. Both the
intersection orientation and the symplectic orientation are constant along the
path $t\mapsto\ker B_{k,t}$, so they agree at $t=1$. Since this holds for every
stratum $D_I$ and every subcollection $A$, all intersections are
$\omega$-symplectic and positively oriented. Thus $D\cup H_1\cup\cdots\cup H_q$
is an $\omega$-symplectic snc divisor. This proves the theorem.

For later use, note that the intersection tangent spaces are uniformly
$O(k^{-1/2})$-close to $J$-complex subspaces. Hence, for any other fixed
symplectic form taming $J$, the same divisors are symplectic snc for all
sufficiently large $k$.
\end{proof}


%------------------------------------------------------------------------------

\section{Regularizations and the local normal model}
\label{SS: Regularizations and local model}

We use the normal-bundle charts of a regularization and their compatibility
under passage between strata. Let $D=D_1\cup\cdots\cup D_r$ be a symplectic snc
divisor in the symplectic $2n$-manifold $(X,\omega)$. For
$I\subseteq\{1,\ldots,r\}$ write $D_I=\bigcap_{i\in I}D_i$, with
$D_\emptyset=X$, and
\[
N_I=N_{D_I|X}=\bigoplus_{i\in I}N_i|_{D_I},
\qquad N_i=N_{D_i|X}.
\]
An $\omega$-regularization $\Psi$ in the sense of \cite{TMZ1} consists, in
particular, of compatible diffeomorphisms
$\Psi_I:\cN_I\to U_I$ from neighborhoods of the zero sections in $N_I$ to
neighborhoods of $D_I$ in $X$. The compatibility conditions identify the
normal-bundle charts associated to nested strata. Moreover,
\cite[Thm.\,2.13]{TMZ1} states that, after deforming $\omega$ through
cohomologous symplectic forms for which $D$ remains a symplectic snc divisor,
such a regularization exists.

We first recall the local symplectic form used in a regularization. Let
$(Y,\omega_Y)$ be a symplectic manifold, let $\pi:N\to Y$ be a symplectic vector
bundle with fiberwise form $\Omega$, and choose a connection on $N$. Writing
$\Omega_\nabla$ for the $2$-form extension of $\Omega$ which vanishes on the
horizontal distribution and $\zeta$ for the fiberwise radial vector field, set
\begin{equation}
\label{Eqn: regularized symplectic form}
\omega_N=\pi^*\omega_Y+\frac12d\bigl(\iota_\zeta\Omega_\nabla\bigr).
\end{equation}
This is symplectic near the zero section and, with respect to the canonical
splitting $T_N|_Y=TY\oplus N$, satisfies
\begin{equation}
\label{Eqn: regularized splitting}
\omega_N|_{T_N|_Y}=\omega_Y\oplus\Omega.
\end{equation}
The regularization maps are symplectomorphisms from forms of this type to the
restriction of $\omega$ to $U_I$.

Now let $J_Y$ be an $\omega_Y$-tame almost complex structure on $Y$, and
suppose that the summands of $N$ have pseudo-holomorphic line-bundle structures
over $(Y,J_Y)$ whose fiberwise complex structures are tamed by $\Omega$. Their
direct sum is a pseudo-holomorphic vector bundle, and the construction in
Section~\ref{SS: Pshol line bundles} gives an almost complex structure
$J_N$ on its total space. Along the zero section,
\[
J_N|_{T_N|_Y}=J_Y\oplus J_{\mathrm{fib}}.
\]
Together with \eqref{Eqn: regularized splitting}, this shows that $J_N$ is
tamed by $\omega_N$ along the zero section, and hence on a neighborhood of it
by openness of tameness.

Applying this construction to $N_I$ and pushing forward by $\Psi_I$ produces
an almost complex structure near $D_I$. In a local pseudo-holomorphic frame of
$N_I=\bigoplus_{i\in I}N_i|_{D_I}$, let $z_i$ be the linear coordinate on the
$i$-th summand. The functions $z_i$ are $J_{N_I}$-holomorphic, and
\[
\Psi_I^{-1}(D_i)=\{z_i=0\},\qquad i\in I.
\]
Thus, provided the normal-bundle structures agree on overlaps of the
regularization charts, the pushed-forward structure makes $D$ a
$J$-holomorphic snc divisor.

We record the compatibility needed below. For $0\le k\le n-1$, let $U(k)$
be the regularized neighborhood of the union of strata of complex dimension
at most $k$, and set
\[
D(k)=D\cap U(k),\qquad D_i(k)=D_i\cap U(k).
\]
An almost complex structure on $D(k)$ means a compatible collection of almost
complex structures on the smooth branches $D_i(k)$ and all their
intersections.

An almost complex structure $J$ on $U(k)$ is
\emph{$\Psi$-regularized} if, for every regularization chart
$\Psi_I:\cN_I\to U_I$ contained in $U(k)$, the pullback $\Psi_I^*J$ is the
total-space almost complex structure defined by $J|_{D_I}$ and
pseudo-holomorphic structures on the summands of
$N_I=\bigoplus_{i\in I}N_i|_{D_I}$.

\begin{lemma}
\label{Lem: regularized local model} Fix a regularization $\Psi$, and set
$U(-1)=\emptyset$. Suppose that a $\Psi$-regularized almost complex structure
has already been constructed on $U(k-1)$. Let $\bar J$ be a compatible almost
complex structure on $D(k)$ whose germ along $D(k-1)$ is induced by this
structure. Suppose moreover that, for every $i$, the complex normal line bundle
$N_i|_{D_i(k)}$ has a $\bar J$-holomorphic structure extending the structure on
$N_i|_{D_i(k-1)}$ induced by the model on $U(k-1)$.

Then, after shrinking the regularization domains, these data determine an
almost complex structure $J$ on $U(k)$ such that $\Psi_I^*J$ is the
total-space almost complex structure defined by $\bar J|_{D_I}$ and the
pseudo-holomorphic structures on the summands of $N_I$. Conversely, every
almost complex structure of this form induces the stated compatible data on
$D(k)$ and the normal line bundles.
\end{lemma}

\begin{proof}
Over a stratum $D_I$, take the direct sum of the prescribed
pseudo-holomorphic structures on $N_i|_{D_I}$, $i\in I$, form the associated
total-space almost complex structure on $N_I$, and push it forward by
$\Psi_I$. If $D_{I'}\subset D_I$, the compatibility maps in the regularization
identify the restriction of this normal model with the model constructed from
$N_{I'}$. On the part lying over $D(k-1)$ this is the compatibility assumed in
the statement. Hence the pushed-forward structures agree on the overlaps of
the charts $\Psi_I$, after their domains are shrunk, and define $J$ on
$U(k)$.

Conversely, pulling such a $J$ back by $\Psi_I$ and restricting to the zero
section gives $\bar J|_{D_I}$. Restriction to each normal summand gives its
pseudo-holomorphic line-bundle structure. Compatibility of these structures
under inclusions of strata follows from the compatibility of the
regularization maps.
\end{proof}

An almost complex structure obtained in this way will be called
\emph{$\Psi$-regularized}. The preceding discussion shows that a
$\Psi$-regularized structure makes $D\cap U(k)$ a $J$-holomorphic snc divisor.
It also shows that if $\bar J$ is tamed by the restrictions of $\omega$ to all
strata and the fiberwise complex structures are tamed by the normal symplectic
forms, then $J$ is $\omega$-tame after shrinking $U(k)$.

We finally recall the standard extension used at the end of the proof.

\begin{lemma}
\label{Lem: Extension of J from divisor}
Let $\Psi$ be an $\omega$-regularization of $D$, and let $J_U$ be a
$\Psi$-regularized almost complex structure on a neighborhood $U$ of $D$.
Assume that its restriction to every stratum is tamed by the corresponding
restriction of $\omega$, and that the fiberwise complex structures on the
normal line bundles are tamed by their normal symplectic forms. After shrinking
$U$, the structure $J_U$ extends to an $\omega$-tame almost complex structure
on $X$.
\end{lemma}

\begin{proof}
By \eqref{Eqn: regularized splitting} and openness of tameness, $J_U$ is
$\omega$-tame on a smaller neighborhood of $D$. The bundle of
$\omega$-tame linear complex structures has nonempty contractible fibers.
Using the convex parametrization of these fibers from
\cite[Prop.\,2.5.13, Proof~3]{McDuffSalamon} and a cutoff function,
one extends $J_U$ from a still smaller neighborhood of $D$ to an
$\omega$-tame almost complex structure on all of $X$.
\end{proof}

%------------------------------------------------------------------------------

\section{Proof of the Main Theorem}
\label{SS: Proof of existence theorem}

\begin{proof}
The construction has two distinct parts. We first enlarge $D$ so that every
normal line on every stratum has a divisor representative avoiding that
stratum. We then use these representatives to construct compatible
pseudo-holomorphic normal-bundle structures by induction on the dimension
of the strata.

Write $X=M$, $\dim_\RR X=2n$, and enumerate the components of $D$ as
$D_1,\ldots,D_r$.

\medskip
\noindent\textsc{Step I}. \emph{Auxiliary divisors.}
For each $i=1,\ldots,r$, choose a complex line bundle $L_i$ with a smooth
section $s_{i0}$ whose transverse, positively oriented zero locus is $D_i$.
The standard tubular-neighborhood construction gives such a pair, with
\[
c_1(L_i)=\operatorname{PD}[D_i].
\]
Put $L_0=\CC_X$ and, for notational purposes, $D_0=\emptyset$ and $s_{00}=1$.

Possibly replacing $\omega$ by a cohomologous symplectic form along a path for
which $D$ remains symplectic snc, we may assume by \cite[Thm.\,2.13]{TMZ1} that
$D$ admits an $\omega$-regularization.

Choose a rational symplectic form $\omega_\QQ$ sufficiently close to
$\omega$. After a sufficiently small regularizing deformation, still denoted
$\omega_\QQ$, the divisor $D$ admits an $\omega_\QQ$-regularization. We may
choose an $\omega_\QQ$-compatible almost complex structure $J_0$ adapted to
this regularization and tamed by $\omega$. Choose $a\in\NN$ such that
$a[\omega_\QQ]$ is integral, and let $P$ be a line bundle with
$c_1(P)=a[\omega_\QQ]$.

Apply Theorem~\ref{Thm: strong Bertini}, using $J_0$ as the auxiliary almost
complex structure, to $n+1$ copies of each of the $r+1$ line bundles
\[
L_0,L_1,\ldots,L_r
\]
on the integral symplectic manifold $(X,a\omega_\QQ)$. For sufficiently large
$N$, it gives, for $i=0,\ldots,r$ and $\mu=1,\ldots,n+1$, sections
$s_{i\mu}$ of
\[
L_{i\mu}:=L_i\otimes P^{\otimes N}
\]
whose zero loci $H_{i\mu}$, together with $D$, form an $a\omega_\QQ$-symplectic
snc divisor. By the final observation in the proof of Theorem~\ref{Thm: strong
Bertini}, after increasing $N$ this divisor is also $\omega$-symplectic snc.
Thus
\begin{equation}
\label{Eqn: tilde D}
\widetilde D
=
D\cup\bigcup_{\substack{0\le i\le r\\1\le\mu\le n+1}}H_{i\mu}
\end{equation}
is an $\omega$-symplectic snc divisor. We refer to the nonempty intersections
$\widetilde D_A$, which may be disconnected, as strata. All dimensions of strata
below are complex.

Set $Q=P^{\otimes N}$ and write $\delta_i=[D_i]$, with $\delta_0=0$, and
$h=\operatorname{PD}(c_1(Q))$. For $\mu>0$, we regard all $L_{i\mu}$ as the
same line bundle $L_i\otimes Q$; the index $\mu$ distinguishes the sections
$s_{i\mu}$ and their zero loci. Thus
\[
[H_{i\mu}]=\delta_i+h,\qquad
L_{i\mu}=L_i\otimes Q.
\]
We include the original components in this notation by setting
\[
H_{i0}=D_i,\qquad L_{i0}=L_i,\qquad 1\le i\le r.
\]
For $\mu,\nu>0$ we then have
\begin{equation}
\label{Eqn: isos of line bundles}
L_{i\mu}=L_{i\nu}\quad(0\le i\le r),
\qquad
L_{i\mu}\otimes L_{0\nu}^{\vee}\simeq L_{i0}
\quad(1\le i\le r),
\end{equation}
where the second identification is induced by
$Q\otimes Q^\vee\simeq\CC_X$.
Let
\[
\Omega=
\{(i,0)\mid 1\le i\le r\}
\cup
\{(i,\mu)\mid 0\le i\le r,\ 1\le\mu\le n+1\}
\]
index the hypersurfaces occurring in \eqref{Eqn: tilde D}. For
$\alpha=(i,\mu)\in\Omega$, write $\widetilde D_\alpha=H_{i\mu}$,
$L_\alpha=L_{i\mu}$ and $s_\alpha=s_{i\mu}$. For $A\subseteq\Omega$, put
\[
\widetilde D_A=\bigcap_{\alpha\in A}\widetilde D_\alpha,
\qquad
\partial\widetilde D_A
=
\bigcup_{\beta\notin A}\widetilde D_{A\cup\{\beta\}}.
\]
Let $N_\alpha=N_{\widetilde D_\alpha|X}$. Transversality of $s_\alpha$
gives an orientation-preserving isomorphism
\begin{equation}
\label{Eqn: N_alpha versus L_alpha}
\nabla s_\alpha:
N_\alpha\xrightarrow{\;\simeq\;}
L_\alpha|_{\widetilde D_\alpha},
\end{equation}
independent of the chosen connection along the zero locus.

The $n+1$ copies ensure the existence of the successors used below. Since
$\dim_\CC X=n$, at most $n$ distinct components of an snc divisor meet.
Consequently, for every nonempty stratum $\widetilde D_A$ and every family
$\{\widetilde D_{(i,1)},\ldots,\widetilde D_{(i,n+1)}\}$, at least one member
has label outside $A$.

\medskip \noindent\textsc{Step II}. \emph{Successor divisors.} We use the
preceding redundancy to represent each normal line bundle $N_\alpha|_{\widetilde
D_A}$ by the restriction of a divisor none of whose components contains
$\widetilde D_A$. The representative changes upon restriction to a deeper
stratum contained in its support.

Fix $A\subseteq\Omega$ with $\widetilde D_A\ne\emptyset$ and
$\alpha=(i,\mu)\in A$. Endow $\{0,\ldots,n+1\}$ with its cyclic order. For
$\alpha=(i,\mu)\in A$, let $\nu_\alpha(A)$ and $\kappa_\alpha(A)$ be the first
positive indices encountered after $\mu$ such that
\[
(i,\nu_\alpha(A))\notin A,
\qquad
(0,\kappa_\alpha(A))\notin A,
\]
respectively. These indices exist because $|A|\le n$.

Define the formal \emph{successor divisor}
\begin{equation}
\label{Eqn: Dagger divisors}
\widetilde D^\dagger_{\alpha,A}
=
\begin{cases}
\widetilde D_{(i,\nu_\alpha(A))}
-\widetilde D_{(0,\kappa_\alpha(A))},&\mu=0,\\[2mm]
\widetilde D_{(i,\nu_\alpha(A))},&\mu>0.
\end{cases}
\end{equation}
It is homologous to $\widetilde D_\alpha$, and every component occurring on the
right-hand side has label outside $A$. Let $L_{\widetilde D^\dagger_{\alpha,A}}$
denote the line bundle obtained from $\widetilde D^\dagger_{\alpha,A}$ by
replacing each $\widetilde D_\beta$ by $L_\beta$. The identifications in
\eqref{Eqn: isos of line bundles} give a canonical isomorphism
\begin{equation}
\label{Eqn: successor line bundle}
L_{\widetilde D^\dagger_{\alpha,A}}\simeq L_\alpha.
\end{equation}
Combining \eqref{Eqn: N_alpha versus L_alpha} and
\eqref{Eqn: successor line bundle} gives an isotopy class of
orientation-preserving isomorphisms
\begin{equation}
\label{Eqn: smooth iso of line bundles}
N_\alpha|_{\widetilde D_A}
\simeq
L_{\widetilde D^\dagger_{\alpha,A}}|_{\widetilde D_A}.
\end{equation}
In Step~III the isotopy class of this isomorphism defines framings of the normal
line bundle and the divisor line bundle associated to $\widetilde
D_A|_{\widetilde D^\dagger_{\alpha,A}}$. We abbreviate $\widetilde
D^\dagger_\alpha =\widetilde D^\dagger_{\alpha,\{\alpha\}}$.

\medskip
\noindent\textsc{Step III}. \emph{Regularization, framings and induction data.}
Possibly replacing $\omega$ once more by a cohomologous symplectic form along
a path for which $\widetilde D$ remains symplectic snc, we may assume by
\cite[Thm.\,2.13]{TMZ1} that $\widetilde D$ admits an
$\omega$-regularization $\Psi=\{\Psi_A\}$.

Below we use the following topological extension statement. A \emph{framing} of
an oriented $\RR^2$-bundle $L$ relative to an oriented $\RR^2$-bundle $L^\Box$
over the same base is an isotopy class of orientation-preserving isomorphisms
$\chi:L^\Box\to L$. If framings of $L$ and $\widehat L$ relative to the same
$L^\Box$ are represented by $\chi$ and $\widehat\chi$, an isomorphism
$\varphi:L\to\widehat L$ is \emph{framed} if $\widehat\chi^{-1}\varphi\chi$ is
isotopic to the identity of $L^\Box$. Now suppose a framed isomorphism
$\varphi:L|_V\to\widehat L|_V$ is given on a neighborhood $V$ of a closed subset
$C\subset X$. After shrinking $V$ around $C$, it extends to a framed isomorphism
over $X$. Indeed, choose an isotopy $a_t$ from the identity of $L^\Box|_V$ to
$\widehat\chi^{-1}\varphi\chi$ and a function $\rho:X\to[0,1]$ supported in $V$
and equal to one near $C$. The automorphism $a_{\rho(x)}$ on $V$, extended by
the identity outside $V$, gives the extension after conjugation by the framings.

For $0\le k\le n-1$, let $U(k)$ be the regularized neighborhood of the union of
strata of $\widetilde D$ of dimension at most $k$, and set $\widetilde
D(k)=\widetilde D\cap U(k)$, $\widetilde D_\alpha(k)=\widetilde D_\alpha\cap
U(k)$. Suppose temporarily that an almost complex structure $\bar J$ has been
constructed on $\widetilde D(k)$ making the union of proper substrata of every
$\widetilde D_A\subset\widetilde D(k)$ into a $\bar J|_{\widetilde
D_A}$-holomorphic snc divisor. For $\alpha\in A$ define
\begin{equation}
\label{Eqn: successor pshol line}
\scrL^\dagger_{\alpha,A}
=
\scrL_{\widetilde D_A}
\bigl(\widetilde D^\dagger_{\alpha,A}|_{\widetilde D_A}\bigr).
\end{equation}
Here we use the notation for integral combinations of divisors introduced in
Section~\ref{SS: Pshol line bundles}.

We frame the normal line bundle and the line bundle in
\eqref{Eqn: successor pshol line} by the same smooth bundle
\[
L^\Box_\alpha=L_\alpha|_{\widetilde D_\alpha}.
\]
The inverse of \eqref{Eqn: N_alpha versus L_alpha} gives a framing
$\chi^\scrN_\alpha:L^\Box_\alpha\to N_\alpha$. On $\widetilde D_A$, the fixed
isomorphism \eqref{Eqn: successor line bundle}, together with the defining
sections of its divisor factors, gives a framing
$\chi^\scrL_{\alpha,A}:L^\Box_\alpha|_{\widetilde D_A}\to
\scrL^\dagger_{\alpha,A}$ of the underlying smooth bundle. Thus the smooth
isomorphism \eqref{Eqn: smooth iso of line bundles} has a distinguished isotopy
class.

The induction will construct pseudo-holomorphic structures $\scrN_{\alpha,k}$ on
$N_\alpha|_{\widetilde D_\alpha(k)}$ and  $L_\alpha^\Box$-framed
pseudo-holomorphic isomorphisms
\begin{equation}
\label{Eqn: phi alpha A}
\varphi_{\alpha,A}:
\scrN_{\alpha,k}|_{\widetilde D_A}
\xrightarrow{\;\simeq\;}
\scrL^\dagger_{\alpha,A},
\qquad \alpha\in A,\quad \dim_\CC\widetilde D_A\le k.
\end{equation}
Their compatibility records the advancement of successors. Let
$B=A\cup\{\beta\}$ and write $\alpha=(i,\mu)$. Put
\[
\beta_+=(i,\nu_\alpha(A))
\]
and, if $\mu=0$, also put
\[
\beta_-=(0,\kappa_\alpha(A)).
\]
Whenever a successor component is encountered, we use the canonical
isomorphism
\[
\scrL_{\widetilde D_A}(\widetilde D_\beta)|_{\widetilde D_B}
\simeq
\scrN_{\beta,k}|_{\widetilde D_B}
\]
from \eqref{Eqn: divisor-normal-bundle}. The compatibility conditions are as
follows.
\begin{enumerate}
\item
If $\beta\ne\beta_+$ and, when $\mu=0$, also $\beta\ne\beta_-$, then
$\scrL^\dagger_{\alpha,A}|_{\widetilde D_B}
=\scrL^\dagger_{\alpha,B}$, and
\[
\varphi_{\alpha,A}|_{\widetilde D_B}
=
\varphi_{\alpha,B}.
\]
\item
If $\mu>0$ and $\beta=\beta_+$, then
\[
\varphi_{\alpha,A}|_{\widetilde D_B}
=
\varphi_{\beta,B}^{-1}\circ\varphi_{\alpha,B}.
\]
\item
If $\mu=0$ and $\beta=\beta_+$, then
\[
\varphi_{\alpha,A}|_{\widetilde D_B}
=
\bigl(\varphi_{\beta,B}^{-1}\otimes\mathrm{id}\bigr)
\circ\varphi_{\alpha,B}.
\]
\item
If $\mu=0$ and $\beta=\beta_-$, then
\[
\varphi_{\alpha,A}|_{\widetilde D_B}
=
\bigl(\mathrm{id}\otimes\varphi_{\beta,B}^{\vee}\bigr)
\circ\varphi_{\alpha,B}.
\]
\end{enumerate}
The framing isomorphisms define the compatible reference system
\[
r_{\alpha,A}=\chi^\scrL_{\alpha,A}\circ
(\chi^\scrN_\alpha|_{\widetilde D_A})^{-1}.
\]
A compatible system is called \emph{framed} if it is isotopic to this system
through compatible systems of smooth bundle isomorphisms.

Thus the isomorphism simply restricts until a current successor is met. Its
divisor line then restricts to its normal line bundle, which is replaced by the
next successor using the already constructed isomorphism for that component. All
identifications come from \eqref{Eqn: isos of line bundles} and \eqref{Eqn:
divisor-normal-bundle}, and preserve the chosen framings.

These rules are compatible under iteration: the successor isomorphisms
cancel telescopically, separately in the two tensor factors when $\mu=0$.
Thus the induced identification on a deeper stratum is independent of the
chain of inclusions.

\medskip
\noindent\textsc{Step IV}. \emph{Induction over the strata.}
We construct, for $k=0,\ldots,n-1$, a $\Psi$-regularized,
$\omega$-tame almost complex structure $J(k)$ on $U(k)$, and denote its
restriction to $\widetilde D(k)$ by $\bar J(k)$. At the same time, we
construct pseudo-holomorphic structures $\scrN_{\alpha,k}$ on the normal
line bundles over $\widetilde D_\alpha(k)$ and a framed compatible system of 
isomorphisms \eqref{Eqn: phi alpha A}. After shrinking the regularization
domains when necessary, all these data extend those constructed at the
previous stage.

For $k=0$, the strata are points, and the required normal-bundle
structures and framed isomorphisms can be chosen pointwise with tame
fiberwise complex structures. Lemma~\ref{Lem: regularized local model}
then produces the corresponding almost complex structure $J(0)$ on
$U(0)$. This starts the induction.

Assume the data have been constructed on $\widetilde D(k-1)$ and let $\widetilde
D_A$ be a $k$-dimensional stratum. Its boundary $\partial\widetilde D_A$ lies in
$\widetilde D(k-1)$. By the local-model Lemma~\ref{Lem: regularized local
model}, the already constructed almost complex structures and pseudo-holomorphic
normal bundles determine an $\omega$-tame regularized almost complex structure
on a neighborhood of $\partial\widetilde D_A$ in $\widetilde D_A$. Extend its
restriction to $T_{\widetilde D_A}$ over all of $\widetilde D_A$, using the
contractibility of the space of $\omega|_{\widetilde D_A}$-tame almost complex
structures. For distinct $k$-dimensional strata $\widetilde D_A$ and $\widetilde
D_{A'}$, the loci $\widetilde D_A\setminus\partial\widetilde D_A$ and
$\widetilde D_{A'}\setminus\partial\widetilde D_{A'}$ are disjoint. The
extensions may therefore be chosen independently away from the prescribed
boundary neighborhoods. They define $\bar J(k)$.

Fix now $\alpha\in A$. On each proper stratum $\widetilde
D_B\subset\partial\widetilde D_A$, the compatibility rule of Step~III expresses
the required restriction of $\varphi_{\alpha,A}$ in terms of the isomorphisms
already constructed on $\widetilde D(k-1)$. By the induction hypothesis, these
prescriptions agree on intersections and define a single $L^\Box_\alpha$-framed
pseudo-holomorphic isomorphism
\begin{equation}
\label{Eqn: boundary phi}
\varphi^\partial_{\alpha,A}:
\scrN_{\alpha,k-1}|_{\partial\widetilde D_A}
\xrightarrow{\;\simeq\;}
\scrL^\dagger_{\alpha,A}|_{\partial\widetilde D_A}.
\end{equation}
We apply the boundary prescriptions to one compatible framing isotopy from the
preceding stage, using it for every $A$. This proves framedness over the whole
boundary $\partial\widetilde D_A$. For every $B\supsetneq A$, the regularization
supplies a retraction from a neighborhood of $\widetilde D_B$ in $\widetilde
D_A$ to $\widetilde D_B$. These retractions are pseudo-holomorphic for the
regularized models. By compatibility of the regularization and the induction
data, the resulting pullbacks agree on overlaps and define an
$L^\Box_\alpha$-framed pseudo-holomorphic isomorphism on a neighborhood of
$\partial\widetilde D_A$ in $\widetilde D_A$.

The two underlying smooth bundles over $\widetilde D_A$ have the framings
defined in Step~III. The framed extension statement therefore extends
$\varphi^\partial_{\alpha,A}$, after shrinking its neighborhood, to a smooth
$L^\Box_\alpha$-framed isomorphism
\[
\varphi_{\alpha,A}:
N_\alpha|_{\widetilde D_A}
\xrightarrow{\;\simeq\;}
\scrL^\dagger_{\alpha,A}
\]
over all of $\widetilde D_A$. Here we extend the boundary framing isotopy to a
smaller neighborhood, with initial map the restriction of $r_{\alpha,A}$ and
final map the isomorphism constructed above, and use this isotopy in the cutoff
construction. Define the pseudo-holomorphic structure on $N_\alpha|_{\widetilde
D_A}$ by pulling back that of $\scrL^\dagger_{\alpha,A}$ through
$\varphi_{\alpha,A}$. Since $\varphi_{\alpha,A}$ is framed, it is
orientation-preserving. The complex orientation of the resulting structure on
the $\RR^2$-bundle $N_\alpha$ therefore agrees with its symplectic orientation;
in real dimension two this implies that the fiberwise complex structure is tamed
by the normal symplectic form.

By construction the pseudo-holomorphic structure on
$N_\alpha|_{\widetilde D_A}$ agrees with $\scrN_{\alpha,k-1}$ near the
boundary, makes $\varphi_{\alpha,A}$ pseudo-holomorphic, and satisfies the
compatibility rule on every proper stratum. Doing this for all
$\alpha\in A$ and all $k$-dimensional strata completes the compatible data on
$\widetilde D(k)$. By Lemma~\ref{Lem: regularized local model}, these data
determine a $\Psi$-regularized, $\omega$-tame almost complex structure
$J(k)$ on $U(k)$. Since they agree with the preceding data near
$\widetilde D(k-1)$, the regularization domains may be shrunk so that $J(k)$
extends $J(k-1)$. This proves the induction step.

At $k=n-1$ we have therefore obtained a $\Psi$-regularized,
$\omega$-tame almost complex structure $J(n-1)$ on $U(n-1)$. Its
restriction $\bar J(n-1)$ to the whole divisor $\widetilde D$ is compatible
with the pseudo-holomorphic structures constructed on all normal line bundles.

\medskip
\noindent\textsc{Step V}. \emph{Extension to $X$ and local coordinates.}
Set $U=U(n-1)$ and $J_U=J(n-1)$. By
Lemma~\ref{Lem: Extension of J from divisor}, after shrinking $U$ the almost
complex structure $J_U$ extends to an $\omega$-tame almost complex
structure $J$ on all of $X$, without changing it on a smaller neighborhood
of $\widetilde D$. In particular, all local pseudo-holomorphic equations
supplied by the regularized model remain unchanged.

It remains to check the coordinate statement. Let $p\in D$, and let
$\widetilde D_{\alpha_1},\ldots,\widetilde D_{\alpha_m}$ be the branches of
$\widetilde D$ through $p$, ordered so that the first $q$ belong to $D$. In
the regularized normal model at their common stratum there are
$J$-holomorphic fiber coordinates $w_1,\ldots,w_m$ with linearly independent
differentials such that
\[
\widetilde D_{\alpha_i}=\{w_i=0\},\qquad i=1,\ldots,m.
\]
The smooth constant-rank theorem completes $(w_1,\ldots,w_m)$ to smooth
complex-valued coordinates $(z_1,\ldots,z_n)$ centered at $p$, with
$z_i=w_i$ for $i\le m$. After shrinking the coordinate neighborhood $V$,
\[
D\cap V=\{z_1\cdots z_q=0\},
\]
and $z_1,\ldots,z_q$ are $J$-holomorphic.

The two deformations of the symplectic form in Steps~I and III concatenate to a
path of cohomologous symplectic forms along which $D$ remains symplectic snc.
Its initial form is the original symplectic form, while its terminal form is the
form currently denoted by $\omega$. Taking this terminal form as $\omega_1$
gives the path in the Main Theorem. The rational form $\omega_\QQ$ is auxiliary
and does not belong to this path.
\end{proof}

%==============================================================================

\begin{thebibliography}{TMZ1}

\bibitem[Au]{Auroux}
D.~Auroux:
\emph{Asymptotically holomorphic families of symplectic submanifolds},
Geom.\ Funct.\ Anal.~\textbf{7} (1997), 971--995.

\bibitem[Do]{Donaldson}
S.~Donaldson:
\emph{Symplectic submanifolds and almost-complex geometry},
J.\ Differential Geom.~\textbf{44} (1996), 666--705.

\bibitem[HP]{HarveyPolking}
R.~Harvey and J.~Polking:
\emph{Removable singularities of solutions of linear partial differential
equations},
Acta Math.~\textbf{125} (1970), 39--56.

\bibitem[MS]{McDuffSalamon}
D.~McDuff and D.~Salamon:
\emph{Introduction to symplectic topology},
3rd ed., Oxford Graduate Texts in Mathematics~27,
Oxford University Press, 2017.

\bibitem[TMZ]{TMZ1}
M.~F.~Tehrani, M.~McLean and A.~Zinger:
\emph{Normal crossings singularities for symplectic topology},
Adv.\ Math.~\textbf{339} (2018), 672--748.

\end{thebibliography}
\end{document}
